\makeatletter
\def\input@path{{./}{../}{../../}{../../../}{../../../../}{preamble/}{../preamble/}{../../preamble/}{../../../preamble/}{../../../../book2026-2027/preamble/}}
\makeatother
\input{preamble_exams.tex}
\input{graphs/highlights.tex}
% \GTenHighlightsfalse % Uncomment to disable instructional highlights.
% Population/model, sample evidence, design threshold, and result roles.
\newcommand{\Pop}[1]{\HighlightStroke{GTenHLParameter}{#1}}
\newcommand{\Obs}[1]{\HighlightStroke{GTenHLValue}{#1}}
\newcommand{\Design}[1]{\HighlightStroke{GTenHLThird}{#1}}
\newcommand{\Result}[1]{\HighlightStroke{GTenHLOutcome}{#1}}

\title{Grade 11 Global Economics\\Term 3 --- Lesson 1\\Home Practice C4: Significance and power}
\author{}\date{}
\hypersetup{pdftitle={Grade 11 Global Economics - Term 3 - Lesson 1 - C4},pdfsubject={Significance and power}}
\begin{document}
\maketitle
\setcounter{tocdepth}{1}
\setcounter{secnumdepth}{2}
\phantomsection\label{toc}
\tableofcontents
\section*{Evaluated criteria}
\begin{description}
\item[\textbf{C4}] Interprets the significance level and power of a hypothesis test, explaining how sample size, effect size, variability, and the selected significance level influence the test’s ability to detect a false null hypothesis.
\end{description}
\textbf{How to practise.} Write each probability with its condition. In every comparison say which factor changes and what stays fixed. Attempt each task before opening its solution link. Use the return arrows to revisit the problem, group, or contents. All datasets and scenarios in this activity are illustrative.
\clearpage
\section{Introduction}\label{sec:introduction}
\subsection*{Highlight legend}
\phantomsection\label{tab:highlight-legend}
\begin{tabularx}{\linewidth}{@{}>{\raggedright\arraybackslash}p{.22\linewidth}>{\raggedright\arraybackslash}p{.25\linewidth}>{\raggedright\arraybackslash}X@{}}
\toprule
\textbf{Symbol or notation}&\textbf{Concept}&\textbf{Brief meaning}\\\midrule
$\Design{\alpha}$ & \Design{Significance} & Rejection under the null. \\
$\Result{1-\beta(\mu_1)}$ & \Result{Power} & Detection at an alternative. \\
$\Pop{\mu_1-\mu_0}$ & \Pop{Effect size} & Departure from benchmark. \\
$\Obs{n},\Pop{\sigma}$ & \Obs{Size and spread} & Precision drivers. \\
\bottomrule
\end{tabularx}
\subsection{Introductory example 1: A specified detection target}\label{ex:intro-one}
A study tests $H_0:\mu=100$ versus $H_a:\mu>100$ for mean invoice cost in USD. Its design has $\Design{\alpha}=\Design{0.05}$ and $\Result{\text{power}}=\Result{0.80}$ specifically at $\mu_1=105$. In repeated valid studies, 5\% reject when $\mu=100$ whereas 80\% reject when $\mu=105$.
\subsection{Introductory example 2: A smaller increase}\label{ex:intro-two}
Under the same test and level, an increase to 102 USD is harder to detect than an increase to 105 USD. Increasing $n$ while holding true mean, spread, and level fixed generally increases power. It does not change the true economic effect.
\section{Notation Reference and Practice Procedure}\label{sec:notation-and-procedure}
\subsection{Notation reference}\label{tab:notation-reference}
\begin{tabularx}{\linewidth}{@{}>{\raggedright\arraybackslash}p{.16\linewidth}>{\raggedright\arraybackslash}p{.18\linewidth}>{\raggedright\arraybackslash}p{.28\linewidth}>{\raggedright\arraybackslash}X@{}}
\toprule
\textbf{Symbol or notation}&\textbf{Concept}&\textbf{Definition}&\textbf{Meaning in the introductory examples}\\\midrule
$\Design{\alpha}$ & \Design{Significance} & $P(\text{reject }H_0\mid H_0\text{ true})$ for an exact test; nominal control can be approximate. & \textbullet\ \hyperref[ex:intro-one]{Example 1}: 0.05 when $\mu=100$. \\
$\Result{1-\beta(\mu_1)}$ & \Result{Power} & $P(\text{reject }H_0\mid\mu=\mu_1)$ at a specified alternative $\mu_1$. & \textbullet\ \hyperref[ex:intro-one]{Example 1}: 0.80 at $\mu_1=105$ USD. \\
$\Pop{\mu_1-\mu_0}$ & \Pop{Effect size} & Raw effect in measurement units; $(\mu_1-\mu_0)/\sigma$ is standardized. & \textbullet\ \hyperref[ex:intro-one]{Example 1}: 5 USD.\par\textbullet\ \hyperref[ex:intro-two]{Example 2}: 2 USD. \\
$\Obs{n},\Pop{\sigma}$ & \Obs{Size and spread} & In a known-spread mean model, $SE=\sigma/\sqrt n$. & \textbullet\ \hyperref[ex:intro-two]{Example 2}: increasing $n$ reduces standard error at fixed spread. \\
\bottomrule
\end{tabularx}
\subsection{Practice procedure}
\begin{enumerate}
\item State parameter, hypotheses, and the specified alternative for power.
\item Interpret significance under the null and power under that alternative.
\item Use $\beta=1-\text{power}$ at that alternative; do not reverse the conditioning.
\item Change one design factor at a time, connecting spread and size to standard error.
\item Distinguish detection from economic importance. Fix the sampling plan before looking at results.
\end{enumerate}
\section{General Worked Problems}\label{sec:general-worked}
\SectionProblemsTableOfContents{\SectionProblemLink{prob:frequency}{Two repeated-study worlds} \\
\SectionProblemLink{prob:design}{Change one factor at a time}}
\subsection{Two repeated-study worlds}\label{prob:frequency}
A valid test of $H_0:\mu=100$ versus $H_a:\mu>100$ has level 0.05 and power 0.80 at $\mu=105$ USD. Imagine 1,000 repetitions under $\mu=100$ and, separately, 1,000 under $\mu=105$.
\QuestionsTasks[sec:general-worked]{{4}/{sol:frequency}/{Find expected decision counts in each world. Explain why they do not give the probability that the null is true after rejection.}}
\subsubsection*{Worked solution}\label{sol:frequency}
Under $\mu=100$, expected rejection count is $1000(\Design{0.05})=50$, with 950 non-rejections. Under $\mu=105$, expected rejection count is $1000(\Result{0.80})=800$, with 200 non-rejections; $\beta(105)=0.20$. These are expected long-run counts, not guaranteed batch outcomes. They condition on different population truths. To find the probability the null is true among rejected studies requires an additional model for how often the different truths occur. Neither $\alpha$ nor a p-value supplies it.
\secInicio[prob:frequency][sec:general-worked]
\subsection{Change one factor at a time}\label{prob:design}
In four separate comparisons for a right-tailed mean test, change only (a) $n=25$ to 100; (b) true increase from 2 to 5 USD; (c) population SD from 20 to 10 USD; (d) $\alpha=0.05$ to 0.01. Hold all other factors fixed within each comparison.
\QuestionsTasks[sec:general-worked]{{4}/{sol:design}/{Explain the effect on power in each case and the trade-off in the final case.}}
\subsubsection*{Worked solution}\label{sol:design}
(a) $SE=\sigma/\sqrt n$ halves, making the fixed effect easier to detect; power rises. (b) The alternative mean moves farther from the null relative to unchanged spread; power rises. (c) Halving $\sigma$ halves $SE$; power rises. (d) The critical value increases from 1.645 to 2.326, so rejection becomes harder and power falls at the fixed alternative, while false rejection under the null decreases. These comparisons retain the same test direction and model. Smaller $\alpha$ does not make the null more likely to be true.
\secInicio[prob:design][sec:general-worked]
\section{Advanced Problems}\label{sec:advanced}
\SectionProblemsTableOfContents{\SectionProblemLink{prob:formula}{Power at two alternatives}}
\subsection{Power at two alternatives}\label{prob:formula}
Independent normal invoice costs have known $\sigma=20$ USD. Plan $n=100$, $H_0:\mu=100$, $H_a:\mu>100$, $\alpha=0.05$. Supplied standard normal references: $z_{0.95}=1.645$, $\Phi(-0.855)\approx0.1963$, $\Phi(0.645)\approx0.7405$. Here $\Phi$ is the standard normal cumulative probability.
\QuestionsTasks[sec:advanced]{{4}/{sol:formula}/{Find the rejection threshold for the sample mean and power at 105 and 102 USD. Explain a practical limitation.}}
\subsubsection*{Worked solution}\label{sol:formula}
The standard error is $20/\sqrt{100}=2$ USD. Reject when $Z=(\bar X-100)/2>1.645$, equivalently $\bar X>100+1.645(2)=103.29$ USD. At alternative $\mu_1$,
\[
\operatorname{power}(\mu_1)=1-\Phi\left(\frac{103.29-\mu_1}{2}\right).
\]
At 105 USD: $1-\Phi(-0.855)\approx0.8037$. At 102 USD: $1-\Phi(0.645)\approx0.2595$. Good detection for a 5 USD increase can coexist with frequent missed 2 USD increases. Power is a probability at a specified truth, not the observed sample's p-value. Economic relevance depends on the procurement budget and response costs. This is prospective planning with fixed $n$; repeatedly extending a sample until significance appears changes false-rejection control.
\secInicio[prob:formula][sec:advanced]
\end{document}

